\documentclass[11pt,a4paper]{article}
\usepackage[spanish,es-noquoting]{babel}
\usepackage[top=2.5cm, bottom=2.5cm, left=2.5cm, right=2.5cm]{geometry}
\usepackage{fontspec}
\usepackage{amsmath,amssymb}
\usepackage{enumitem}
\usepackage{microtype}

\setmainfont{Calibri}
\linespread{1.15}

\begin{document}
\pagestyle{empty}

\noindent\textbf{1.} En la superficie de un exoplaneta cuya aceleración de la gravedad es $g = 18\text{ m/s}^2$, se efectúa el lanzamiento de un objeto explorador. La ecuación vectorial de posición que describe su movimiento viene determinada por:
\[
\vec{r}(t) = 
\begin{cases}
x(t) = 24t + 5 \\
y(t) = -9t^2 + 25t + 10
\end{cases}
\quad (\text{con } x, y \text{ en metros y } t \text{ en segundos})
\]

\noindent Determinar:

\begin{enumerate}[label=\textbf{\alph*)}, leftmargin=2em, itemsep=6pt]
    \item El vector de posición en el instante inicial.
    \item El vector desplazamiento en los dos primeros segundos.
    \item El vector velocidad media en los dos primeros segundos y su módulo.
    \item La ecuación de la trayectoria en función de la coordenada horizontal, $y(x)$.
    \item El vector velocidad y su módulo en el instante $t = 1\text{ s}$.
    \item El vector aceleración y su módulo en ese instante ($t = 1\text{ s}$).
    \item La aceleración tangencial ($a_t$) en ese instante.
    \item La aceleración normal ($a_n$) en ese instante.
    \item El radio de curvatura ($\rho$) de la trayectoria en ese instante.
\end{enumerate}

\end{document}
